\sscp{Issues with transcription}{02-04-05}
There are many issues relating to problems that arise from the
transcription of words that muddy the picture for comparative-historical purposes.
We discuss here only a sample for illustrative purposes.

Are intervocalic semi-vowels phonemic? phonetic? or simply transcription
variants? For example, words like \it{sawan, awan,
dawe, tewa} are fairly straightforward.
But what about \it{duwe {\tl} due,
ua {\tl} uwa}, \it{leya {\tl} leiya {\tl} lea}? Or \it{tewa {\tl}
teiwa {\tl} taiwa}? This issue is relevant to many languages
in the region, and to understanding retentions (or not) of PMP *w
and *y, yet many investigators do not address them.

There are wide discrepancies of transcription in distinguishing
true vowel-initial words from glottal stop-initial words,
often making comparative work difficult,
and sometimes under-representing \tsc{1sg} or nominal prefixes.
Is it /eu/ ‘I go’? Or /ʔ-eu/ ‘\tsc{1sg}-go’? Is it /ai/
and /au/ ‘tree, wood’? Or /ʔai/
and /ʔau/ (reflecting the historical *k of PMP *kahiw)?
And how do we know? \citet[60--62]{Tamelan2021} provides several tests
for the \ili{Dela} language to determine whether a glottal is phonemic or not.
There are some languages in which PMP *k > \it{ʔ},
but it might not be written in any position by the investigators,
or is written only between vowels but not initially or finally,
or is written only irregularly although it actually occurs regularly.
We have encountered all of these issues.
There are other languages in which PMP *k > {\0} in all positions.
But the distinction between *k > \it{ʔ}
and *k > {\0} can be significant for comparative purposes and for subgrouping.
These examples also illustrate that orthographies used for linguistic
resources should perhaps be different on some points from practical
orthographies for popular consumption.

Several older Dutch and German sources in the region indicate
an intervocalic glottal with a diaeresis over the second vowel.
So orthographic ‹seï› for /seʔi/,
‹haü› for /haʔu/, ‹foeöe› for /fuʔu/,
‹djaä› for /ʤaʔa/, etc. (see \citealp{Grimes1999}).
Yet in other sources the diaeresis may represent that the
two vowels are articulated as separate syllables,
rather than coalescing as a phonetic diphthong.

Given so much variation in transcription,
consequently in several Holle lists it is not surprising that
the introductory notes indicate a diaeresis over a vowel
indicates an intervocalic glottal.
Yet this is also said for languages that have no intervocalic
glottals (such as \ili{Buru}
and others in central Maluku).
Diacritics of various sorts (with little consistency across languages),
are used to indicate: glottal,
stressed syllable, unstressed syllable,
length, schwa, close vowel, open vowel,
cancelling Dutch orthographic practices (i.e.
a diacritic may indicate ‹oe› should be read as a vowel
sequence /oe/, not as /u/), and other things.
In different lists (and sometimes in the same list),
‹a›, ‹á›, ‹à›, ‹ä›,
‹â›, ‹å›, ‹ā›, ‹a›, and ‹aa›,
could indicate the same thing,
or they could indicate different things.
One must know something about the language to decode the diacritics,
and phonemic differences need to be retained in the documentation
and reflected in comparative-historical discussion.

Many Austronesian languages in the target region have VV sequences of like vowels.
These are phonetically longer than a single V,
and measured instrumentally average around 20--45{\%} longer
in duration than a single V
(see \citealp[95--99]{Edwards2020}
and \citealp{Grimes2024}).
The duration of VV sequences of like
and unlike vowels is often only trivially different,
with unlike sequences tending to be slightly longer.
Phonemically these double VV are often analysed as two syllable
peaks, and this is reflected in the patterns of word stress,
affixation, compounding, and reduplication \citep{CulhaneEdwards2021}.\footnote{
		Unfortunately this very common
		and significant feature of the vowel systems of the region is
		only mentioned in passing by \citet{Hajek2010}.}

Native speakers may be inconsistent in writing the difference
between double vowels and single vowels, as their main orthographic model is usually
Indonesian -- a language which does not distinguish double vowels.
However, when they become aware of the difference,
they consistently write the differences between a single
and double vowel in many languages in the region.
Some same VV sequences arise from the loss of a historical
consonant as illustrated in \trf{02-07}
(STB = \tsc{Seram-Tanimbar-Bomberai}).

\begin{table}
	\centering\tcp{VV sequences from loss of historical *C (consonant)}{02-07}
	\st{4.25pt}\begin{tabular}{ll|l@{ }l|l@{ }l|l}\lsptoprule
&&&&&&trunk, base,\\
&local gloss&\mc{2}{l|}{sail}&\mc{2}{l|}{name}&source\\
Subgroup&PMP&\mc{2}{l|}{*layaR}&\mc{2}{l|}{*ŋajan}&*puqun\\ \midrule
\tsc{\ili{Bima}-Lembata}&\ili{Dhao}&\mc{2}{l|}{{\hr}lai}&\mc{2}{l|}{{\hr}ŋara}&\tcg{(kapua)}\\
\tsc{Timor-Babar}&\ili{Tii}&\mc{2}{l|}{{\hr}laa}&\mc{2}{l|}{{\hr}nade}&{\hr}huu\\
\tsc{Timor-Babar}&\ili{Tetun}&\mc{2}{l|}{{\hr}laa-n}&\mc{2}{l|}{{\hr}nara-n}&{\hr}huun\\
\tsc{\ili{Central} Timor}&\ili{Kemak}&\mc{2}{l|}{}&\mc{2}{l|}{{\hr}galan}&{\hr}puun\\
\tsc{Sula-Buru}&\ili{Buru}&{\hr}laa &‘sail’&{\hr}ŋaa-n &‘name’&{\hr}puun\\
&&{\hr}laa-n&‘dorsal fin’&{\hr}ŋaa-t&‘rank’&\\
STB&\ili{Fordata}&\mc{2}{l|}{{\hr}laar}&\mc{2}{l|}{{\hr}nara-n}&\\
STB&\ili{Geser}&\mc{2}{l|}{{\hr}laar}&\mc{2}{l|}{{\hr}ŋasa-n}&\\
		\lspbottomrule
	\end{tabular}
\end{table}

It is also the case that long vowels (or disyllabic double vowels)
can be indicated in the primary source material,
but not preserved in the comparative material.
For example, \citet{Jonker1908,Jonker1915}
indicated some long vowels with an acute accent,
others with a grave accent,
and still others not at all.
But even where he did indicate them,
this important information may not be preserved in comparative
works that use him as a source
(see discussion in \citeauthor{Edwards2018a} [\citeyear[366]{Edwards2018a}]
and \citeauthor{Edwards2021} [\citeyear[8--12]{Edwards2021}]).

\citet{BlustTrussel2020} under PCEMP *maya ‘tongue’ list ``Roti''
‹ma-k› and ``Atoni'' ‹ma-f›.
Our own notes, as well as published native-authored materials
for several varieties of Rote,
and several varieties of \ili{Meto} have \it{maa-k}
and \it{maa-f} respectively
(and the same root with other genitive enclitics as well).
*y is regularly lost,
but both historical vowels are retained.
How these words are transcribed can help reveal the historical
patterns, or mask them and contribute to the confusion.

\citet[19-21]{vanKlinken1999} indicates double vowels in \ili{Tetun},
which she analysed phonemically as two syllables,
with an acute accent over the vowel,
as in ‹hare› ‘rice’,
‹haré› ‘see’, which is analysed phonemically as /haree/.
For comparative purposes these need to be written as VV sequences of like vowels.
For example, orthographic ‹tún› ‘descend’, needs to be reinterpreted
as \it{tuun} in the comparative literature to track its relationship
with PMP *tuRun ‘descend’,
with the regular loss of *R,
but no loss of a historical vowel. Data sources vary.
When Grimes was driving from Dili to \ili{Kupang} a road sign read
‹hatun velocidade› intending /ha-tuun/ ‘cause to descend = decrease (speed)’.
Without knowledge of the language ‹hatun› could be misinterpreted
as /hatun/, which looks like a reflex of PMP *batu ‘stone, rock’.

In the \ili{Asilulu} dictionary,
\citet{Collins2003a} explains “The colon (:) marks long vowels,
which contrast phonemically with short vowels.” (p.
xxvii) “perhaps a more consistent orthography would have represented
these long vowels as double vowels (\it{aa},
\it{ii}, and so forth).” (p. xxxvi).
It is also the case that for a single language,
one investigator may indicate the long/double vowels,
while another investigator may be oblivious to their existence.

Informed clarifications are extremely helpful,
such as the information found in \citet{Nivens2017} that in the
comparative \ili{Aru} data, transcriptions are adapted to IPA standards,
with the following exceptions: the letter ‹y› is used for
the palatal glide [j],
\tsc{Tarangan} mid-close [e] and [o] are transcribed ‹ê›
and ‹ô›; mid-open [ɛ]
and [ɔ] are transcribed ‹e›
and ‹o›; and unexpected stress is marked by an acute accent (the
normal pattern is penultimate stress,
or final stress if the final syllable contains a mid-vowel).
The symbol ‹ʍ› represents a voiceless labiovelar fricative.
Such clarification further illustrates our challenges for unravelling
the wide variations of what is indicated throughout the region
by various diacritics over vowels.
Unfortunately, many sources of data provide no information about
the transcription practices employed.

In unravelling all this confusion,
we have often had to triangulate with multiple sources,
including consulting more modern sources,
more reliable sources on neighbouring languages,
consulting with scholars who have worked on those languages to
clarify what has been said by them or by others,
and occasionally we have found native speakers or gone to the
region to better understand the phenomena hinted at in the literature
or transcription of those languages.


